\section{Data Benchmarking and Preprocessing}
\label{sec:data_benchmark}

There are many open-source pretraining datasets across various data domains, especially for English, Code, and Math~\cite{olmo20252olmo2furious,smol2,yiwen-etal-2025-yulan,DCLM,su2025nemotroncctransformingcommoncrawl,stack-v2,zhou2025megamath}.
However, constructing a high-quality pretraining corpus remains a non-trivial task due to two primary challenges. First, it is challenging to measure the quality of diverse datasets and determine the optimal strategy for selecting and mixing data from heterogeneous sources.
Second, preprocessing these datasets is both resource-intensive and technically complex.
Given the large scale of pretraining data and complex operations like deduplication, the preprocessing pipeline incurs substantial computational overhead and engineering complexity.

To mitigate these issues, (1) we propose to benchmark primary datasets (e.g., DCLM-Baseline~\cite{DCLM}, Fineweb-Edu~\cite{fineweb}) by quantiles of quality scores. We train reference models over the data subsets around a series of quantiles of quality scores, and then analyze how the resulting benchmark performance varies with data distribution, which is reflected by quality scores. (2) We develop a user-friendly Spark-based data preprocessing framework to efficiently process large-scale pretraining datasets. Moreover, we exploit the Chukonu~\cite{chukonu} framework to reduce the preprocess overhead. These explorations on data dimensions lay a solid foundation for our training and future work.

\subsection{Benchmarking Dataset By Quality-Score Quantiles}\label{sec:data_benchmarking_subsec}

\begin{figure}[ht]
    \centering
    \includegraphics[width=0.85\linewidth]{quantile_benchmarking_diagram.pdf}
    \caption{\textbf{Illustration of the Quantile Benchmarking Process.}
    (1) Given a series of target quantiles (e.g., 0\%, 20\%, 40\%, 60\%, 80\%), we select a data chunk around each target quantile as a probing dataset. 
    (2) A small-scale reference model is then evaluated on each probing dataset under two settings: training from scratch, or resuming from a checkpoint for continual training.}
    \label{fig:quantile_benchmark_process}
\end{figure}

\paragraph{Background and Motivation.} Most open-source datasets have been through a preprocessing pipeline, which primarily incorporates steps of quality scoring and data filtering by score. These score labels are typically released for these open-source datasets. Therefore, we can select a subset hypothesized to be of high quality based on sample quality scores. However, samples between different datasets are hard to compare, considering heterogeneous quality metrics. When scorers are available, it is possible to score both datasets. But as more datasets and quality metrics are included, it is hard to scale up and judge by multiple quality metrics. DCLM~\cite{DCLM} proposes to benchmark datasets or quality scores by filtering and feeding datasets into a standard series of models across scales. However, when facing a practical pretraining setup, the top-k filtering and multi-scale benchmarking can be challenging: the cost of full benchmarking is prohibitive, and we need to ablate over different filtering ratios to balance quality and quantity of the filtered dataset. 

\subsubsection{Quantile Benchmarking Methodology}

\textit{Quantile Data Benchmarking} refers to the systematic evaluation of dataset quality by training reference models on data subsets stratified by quality score percentiles, enabling empirical comparison of dataset characteristics across different quality ranges. Instead of relying solely on top-$k$ filtering, we design a small-scale evaluation process across a range of quality score quantiles to benchmark dataset quality. In practice, preprocessed open-source datasets typically provide a quality metric for each data sample. These quality scores reflect specific characteristics of the data and can vary significantly. It is commonly expected that using higher-quality data samples in training should lead to better model performance. Motivated by this intuition, we propose a straightforward approach, which has not been systematically explored in prior work.

For a target dataset, we first determine a series of quality score quantiles, such as top 0\%, 10\%, 20\%, $\dots$, 80\%. At each quantile, we select a fixed-size subset of the dataset. In our implementation, starting from the data sample ranked at the top 10\% in terms of quality score, we expand the subset by including the next 10B tokens of lower-quality samples to form the probing dataset. We then train a small-scale reference model (e.g., 0.5B parameters) on each of these probing datasets. Finally, we evaluate the resulting models on a set of target benchmarks to record their performances. We refer to this process of evaluating datasets across different quality quantiles as \textit{quantile benchmarking}.

Given the computational cost of training multiple reference models on different probing datasets, we typically apply quantile benchmarking to dominant datasets, such as DCLM-Baseline~\cite{DCLM} and FineWeb-Edu~\cite{fineweb} in the English domain, and FineWeb-Edu-Chinese-V2.1~\cite{fineweb-edu-chinese} in the Chinese domain. Moreover, we measure the utility of each probing dataset in two scenarios: (1) training the reference model from scratch, and (2) continuing training from pretrained checkpoints. Evaluating both scenarios provides a more comprehensive understanding of the target dataset. \Cref{fig:quantile_benchmark_process} illustrates the overall quantile benchmarking process, including quantile data selection and benchmarking experiments under both scenarios. 

While quantile benchmarking introduces additional computational overhead, this cost is marginal compared to the resources required for full pretraining.
For instance, as illustrated in \Cref{fig:quantile_benchmark_results}, we benchmark the DCLM-Baseline across 5 quantiles, ranging from 0\% to 60\% at 15\% intervals.
Although the entire deduplicated DCLM-Baseline comprises 609B tokens (as reported in \Cref{tab:phase1_stats}), our benchmarking process consumes only 8.4B tokens per run, totaling 42B tokens (approximately 6.9\% of the dataset).
Furthermore, because we utilize a reference model of at most 0.6B parameters, the total computational budget constitutes only 2\% of the cost required to train a 2B target model over the full DCLM-Baseline in a single pass, and represents less than 0.6\% of the total pretraining budget of \sys (Explained in \Cref{sec:referece_quantile_benchmarking}).
Given that quantile benchmarking is performed only once per target dataset to guide all subsequent selection decisions, the actual amortized cost is negligible.
Beyond efficiency, these experiments offer significantly higher granularity than training a reference model over a subset via top-$k$ filtering, which merely reflects the average quality of data above a specific threshold rather than the distribution within specific bands.
Consequently, we strategically target this benchmarking at leading datasets, such as DCLM-Baseline, Fineweb-Edu, and Fineweb-Edu-Chinese-V2.1, where the approach exhibits the most favorable benefit-cost ratio.


\subsubsection{Benchmarking Results and Analysis}
Based on our quantile experiment results, we compare models trained on different dataset partitions across various benchmarks. These quantile-based comparisons provide deeper insights into the characteristics of target datasets and offer guidance for data selection and mixing strategies.

\begin{figure}
    \begin{center}
        \begin{subfigure}[b]{0.48\linewidth}
            \centering
            \includegraphics[width=\linewidth]{benchmark_plots_resume/MMLU_resume.pdf}
        \end{subfigure}
        \begin{subfigure}[b]{0.48\linewidth}
            \centering
            \includegraphics[width=\linewidth]{benchmark_plots_resume/WinoGrande_resume.pdf}
        \end{subfigure}
    \end{center}
    \caption{Representative results showing task-dependent dataset characteristics: FineWeb-Edu excels on know\-ledge-intensive tasks (MMLU) while DCLM-Baseline performs better on commonsense reasoning (WinoGrande).}
    \label{fig:quantile_benchmark_results}
\end{figure}

As an illustrative example, we present quantile experiments on both Fineweb-Edu and DCLM-Baseline, offering a complementary perspective to previous analyses~\cite{su2025nemotroncctransformingcommoncrawl,wettig2025organizewebconstructingdomains}. The representative comparison sees \Cref{fig:quantile_benchmark_results} and full comparison results can refer to \Cref{fig:quantile_benchmark_understanding,fig:quantile_benchmarks_knowledge}. Our investigation aims to deepen the understanding of these representative open-source datasets and identify their key differences and commonalities, which we summarize as follows:

\begin{enumerate}
    \item[(1)] \textbf{Task-dependent dataset superiority.} Fineweb-Edu generally demonstrates superior performance on academic and encyclopedic benchmarks, including MMLU~\cite{mmlu} and Common Sense QA (CSQA)~\cite{commonsenseqa}, as well as reading comprehension tasks like BoolQ~\cite{clark-etal-2019-boolq}. In contrast, DCLM-Baseline exhibits slight advantages on situated commonsense reasoning, such as PIQA~\cite{piqa}, Social IQa~\cite{socialiqa}, HellaSwag~\cite{hellaswag}, and WinoGrande~\cite{winogrande}. This divergence suggests that FineWeb-Edu excels in tasks requiring more structural knowledge and formal semantics, while DCLM may benefit tasks relying on intuitive scenario-based reasoning. Representative comparisons are illustrated in \Cref{fig:quantile_benchmark_results} where Fineweb-Edu induces better results on MMLU while DCLM-Baseline outperforms on WinoGrande. More comprehensive results are presented in \Cref{fig:quantile_benchmarks_knowledge,fig:quantile_benchmark_understanding}, which respectively highlight academic knowledge and formal reasoning, and situated commonsense reasoning.
    \item[(2)] \textbf{Substantial within-dataset heterogeneity.} Data quality varies considerably within individual datasets. For instance, in the continual training (resume) scenario, DCLM-Baseline exhibits a 2\% performance difference on MMLU between the top 0\% and top 60\% quantiles, while showing an even more pronounced 8\% variation on ARC-Easy across the same quality range, as shown in \Cref{tab:reasoning_knowledge}. This substantial heterogeneity underscores the importance of quality-aware data selection and training strategies.
    
    \item[(3)] \textbf{Consistency across training scenarios.} The relative superiority relationships between datasets remain largely consistent across both continual training (resume) and from-scratch (run) scenarios, as demonstrated in \Cref{fig:quantile_benchmark_understanding,fig:quantile_benchmarks_knowledge}. However, we observe occasional deviations in specific quantile ranges, suggesting that training dynamics may influence relative dataset effectiveness.
    
    \item[(4)] \textbf{Non-monotonic quality-performance relationships.} Benchmark performance does not necessarily increase monotonically with quality scores. As shown in \Cref{fig:quantile_benchmark_understanding,fig:quantile_benchmarks_knowledge}, increasing quality scores, measured by the fineweb-edu classifier for Fineweb-Edu and the FastText score for DCLM-Baseline, can paradoxically lead to decreased performance on HellaSwag and PIQA. This finding calls into question the universal applicability of quality metrics employed in current leading open-source datasets, and highlights the task-specific nature of data quality assessment.
\end{enumerate}

In summary, our quantile experiments reveal that (i) datasets exhibit substantial internal heterogeneity, and (ii) the relative superiority of both datasets and their quality partitions is highly dependent on the target capability of interest. 
Quality assessment is task-dependent and context-specific, necessitating careful consideration when comparing datasets across different capability dimensions. More details of implementation and discussions are presented in \Cref{sec:referece_quantile_benchmarking}.

\subsubsection{Implications for Data Selection}
These findings inform our data mixing and training strategies in the following ways:

\begin{enumerate}
    \item[(1)] \textbf{Curriculum learning with selective repetition.} Beyond the conventional practice of filtering low-quality data, we propose strategically scheduling high-quality data partitions toward later training stages (curriculum learning) while applying higher repetition rates to these partitions compared to lower-quality data (selective multi-epoch training). This approach leverages within-dataset quality variation to enhance training efficiency (details in \Cref{sec:multi_phase_curriculum}).
    
    \item[(2)] \textbf{Benchmark-guided mixing ratios.} Given a representative benchmark aligned with a target capability, such as MMLU for knowledge-intensive tasks, quantile comparisons can guide inter-dataset mixing ratios. For example, as illustrated in the left panel of \Cref{fig:quantile_benchmark_results}, the entire Fineweb-Edu dataset exhibits performance roughly comparable to the top 30\% partition of DCLM-Baseline on MMLU, suggesting appropriate relative sampling rates for knowledge-focused pretraining. 
    Concretely, in Phase 2 (\Cref{tab:phase2_stats}), we use the whole Fineweb-Edu dataset while using only the top 33.4\% DCLM-Baseline dataset. In the latter phase, the relative ratio of DCLM-Baseline is further pulled down, shown in \Cref{tab:phase3_stats,tab:phase4_stats,tab:phase5_stats}. The heuristic mixing ratio design lays a foundation for more strategies introduced in \Cref{sec:multi_phase_curriculum}.
\end{enumerate}

We acknowledge that the current analysis remains primarily qualitative and coarse-grained. More fine-grained, quantitative frameworks for dataset comparison and mixing ratio optimization represent promising directions for future research.

\subsection{Data Processing Framework}
\label{sec:data_preprocessing_framework}

To address the challenges of data processing, our data processing framework is designed to satisfy three critical requirements:
\begin{enumerate}
    \item \textbf{Reproducibility:} Given that the training dataset of \sys is composed of various open-source datasets, the framework should be able to reconstruct the exact dataset from these original sources with a configuration file.
    \item \textbf{Usability and Scalability:} The framework should support various operations like filtering, deduplication and mixing. Furthermore, this framework should scale to large clusters without additional engineer efforts.
    \item \textbf{High Performance:} To handle hundreds of terabytes of data, the framework must be optimized to reduce computation overhead.
\end{enumerate}

To meet these demands, we developed \textbf{\dataframework}, a distributed data processing framework built on Spark~\cite{zaharia2012spark}.
\dataframework adopts a tree-structured processing pipeline design.
The leaf nodes represent the raw open-source datasets, while internal nodes represent processing operators like filters and samplers.
The root node generates the final mixed training dataset.
With this design, the entire processing pipeline, including dataset sources and operator parameters, can be defined with a YAML configuration file.
This ensures strict reproducibility, enabling researchers to reconstruct the exact training corpus from raw datasets simply by applying the configuration.

As \dataframework is built on Spark, it inherits the programming flexibility and scalability.
We utilize the powerful Spark RDD API to develop complex data processing operators, and rely on the Spark Engine for distributed processing, resource management, and fault tolerance.
This design allows \dataframework to process over 100 TB of data across large-scale clusters with minimal engineering efforts.

Despite Spark's scalability, the overhead of JVM-based execution can become a bottleneck for compute-intensive tasks. To address this, we integrated the Chukonu~\cite{chukonu} framework, utilizing its C++ interface to refactor certain performance-critical operators.
By conducting computations with native C++, we accelerates the processing procedure.
For instance, the optimized MinHash deduplication operator is approximately $2.5\times$ faster than the Spark implementation.





